\section{Evaluation: 50-State Compatibility Census}
\label{sec:evaluation}

We apply the \textsc{SpineCompatibility} score to all 50 US states using the
2020 apportionment congressional seat counts and the 2020 state legislative
chamber sizes from the National Conference of State Legislatures (NCSL).

\subsection{Full Compatibility Table}

Table~\ref{tab:compat} reports all 50 states sorted by compatibility score
$\sigma$ (ascending).  States are classified into three groups:
\emph{Strictly compatible} ($\sigma = 0$, $g = \min(C,S,H)$),
\emph{Partially compatible} ($0 < \sigma < 50$, strong trunk), and
\emph{Structurally incompatible} ($\sigma \geq 50$, weak trunk).

\begin{longtable}{llrrrrrl}
\caption{NestSection compatibility scores for all 50 states (2020 apportionment).
$C$ = congressional seats, $S$ = state senate seats, $H$ = state house seats,
$g = \gcd(C,S,H)$, $\sigma$ = spine compatibility score.}
\label{tab:compat} \\
\toprule
State & Code & $C$ & $S$ & $H$ & $g$ & $\sigma$ & Class \\
\midrule
\endfirsthead
\multicolumn{8}{c}{\tablename\ \thetable\ (continued)} \\
\toprule
State & Code & $C$ & $S$ & $H$ & $g$ & $\sigma$ & Class \\
\midrule
\endhead
\bottomrule
\endfoot
\bottomrule
\endlastfoot
%% STRICTLY COMPATIBLE (sigma = 0) -- 11 states
Alaska        & AK &  1 &  20 &  40 & 1 &   0.00 & Strict \\
Alabama       & AL &  7 &  35 & 105 & 7 &   0.00 & Strict \\
Delaware      & DE &  1 &  21 &  41 & 1 &   0.00 & Strict \\
Montana       & MT &  2 &  50 & 100 & 2 &   0.00 & Strict \\
North Dakota  & ND &  1 &  47 &  94 & 1 &   0.00 & Strict \\
New Hampshire & NH &  2 &  24 & 400 & 2 &   0.00 & Strict \\
Oregon        & OR &  6 &  30 &  60 & 6 &   0.00 & Strict \\
South Dakota  & SD &  1 &  35 &  70 & 1 &   0.00 & Strict \\
Vermont       & VT &  1 &  30 & 150 & 1 &   0.00 & Strict \\
West Virginia & WV &  2 &  34 & 100 & 2 &   0.00 & Strict \\
Wyoming       & WY &  1 &  30 &  60 & 1 &   0.00 & Strict \\
\midrule
%% INCOMPATIBLE (sigma >= 50) -- 39 states
Hawaii        & HI &  2 &  25 &  51 & 1 &  50.00 & Incompatible \\
Iowa          & IA &  4 &  50 & 100 & 2 &  50.00 & Incompatible \\
Idaho         & ID &  2 &  35 &  70 & 1 &  50.00 & Incompatible \\
Louisiana     & LA &  6 &  39 & 105 & 3 &  50.00 & Incompatible \\
Maine         & ME &  2 &  35 & 151 & 1 &  50.00 & Incompatible \\
Mississippi   & MS &  4 &  52 & 122 & 2 &  50.00 & Incompatible \\
Rhode Island  & RI &  2 &  38 &  75 & 1 &  50.00 & Incompatible \\
Arizona       & AZ &  9 &  30 &  60 & 3 &  66.67 & Incompatible \\
Kentucky      & KY &  6 &  38 & 100 & 2 &  66.67 & Incompatible \\
Nebraska      & NE &  3 &  49 &  49 & 1 &  66.67 & Incompatible \\
New Jersey    & NJ & 12 &  40 &  80 & 4 &  66.67 & Incompatible \\
New Mexico    & NM &  3 &  42 &  70 & 1 &  66.67 & Incompatible \\
Tennessee     & TN &  9 &  33 &  99 & 3 &  66.67 & Incompatible \\
Arkansas      & AR &  4 &  35 & 100 & 1 &  75.00 & Incompatible \\
Kansas        & KS &  4 &  40 & 125 & 1 &  75.00 & Incompatible \\
Nevada        & NV &  4 &  21 &  42 & 1 &  75.00 & Incompatible \\
Utah          & UT &  4 &  29 &  75 & 1 &  75.00 & Incompatible \\
Connecticut   & CT &  5 &  36 & 151 & 1 &  80.00 & Incompatible \\
Ohio          & OH & 15 &  33 &  99 & 3 &  80.00 & Incompatible \\
Oklahoma      & OK &  5 &  48 & 101 & 1 &  80.00 & Incompatible \\
Florida       & FL & 28 &  40 & 120 & 4 &  85.71 & Incompatible \\
Georgia       & GA & 14 &  56 & 180 & 2 &  85.71 & Incompatible \\
North Carolina& NC & 14 &  50 & 120 & 2 &  85.71 & Incompatible \\
South Carolina& SC &  7 &  46 & 124 & 1 &  85.71 & Incompatible \\
Colorado      & CO &  8 &  35 &  65 & 1 &  87.50 & Incompatible \\
Maryland      & MD &  8 &  47 & 141 & 1 &  87.50 & Incompatible \\
Minnesota     & MN &  8 &  67 & 134 & 1 &  87.50 & Incompatible \\
Missouri      & MO &  8 &  34 & 163 & 1 &  87.50 & Incompatible \\
Wisconsin     & WI &  8 &  33 &  99 & 1 &  87.50 & Incompatible \\
Indiana       & IN &  9 &  50 & 100 & 1 &  88.89 & Incompatible \\
Massachusetts & MA &  9 &  40 & 160 & 1 &  88.89 & Incompatible \\
California    & CA & 52 &  40 &  80 & 4 &  90.00 & Incompatible \\
Washington    & WA & 10 &  49 &  98 & 1 &  90.00 & Incompatible \\
Virginia      & VA & 11 &  40 & 100 & 1 &  90.91 & Incompatible \\
Michigan      & MI & 13 &  38 & 110 & 1 &  92.31 & Incompatible \\
Illinois      & IL & 17 &  59 & 118 & 1 &  94.12 & Incompatible \\
Pennsylvania  & PA & 17 &  50 & 203 & 1 &  94.12 & Incompatible \\
New York      & NY & 26 &  63 & 150 & 1 &  96.15 & Incompatible \\
Texas         & TX & 38 &  31 & 150 & 1 &  96.77 & Incompatible \\
\end{longtable}

\paragraph{Summary.}  Of the 50 states, 11 are strictly compatible ($\sigma = 0$),
0 fall in the partial range $0 < \sigma < 50$, and 39 are structurally
incompatible ($\sigma \geq 50$).  No state falls in the intermediate partial
range; the distribution is bimodal, clustering at $\sigma = 0$ (single-seat
and small-seat states) and $\sigma \geq 50$ (the large-chamber mismatch regime).

\subsection{Why No Partially Compatible States?}

The gap between strict ($\sigma = 0$) and incompatible ($\sigma \geq 50$)
reflects the structure of small-seat state apportionments.  For $\sigma$ to
fall strictly between 0 and 50, we need $\gcd(C,S,H)/\min(C,S,H) > 1/2$,
i.e.\ $g > m/2$.  This requires a GCD larger than half the minimum seat count.
Among the 11 single-seat congressional states ($C = 1$), $g = 1 = \min$
so $\sigma = 0$ exactly.  For states with $C \geq 2$, achieving $g > C/2$
requires $S$ and $H$ to share more than half of $C$'s prime factors, which
coincides with $g = C$ (strict) in all 50 current apportionments.

\subsection{Case Study 1: Oregon --- Perfect Nesting}

Oregon has $C = 6$, $S = 30$, $H = 60$.  Then:
\begin{align*}
  g &= \gcd(6, 30, 60) = 6 \\
  \sigma &= (1 - 6/6) \times 100 = 0\%
\end{align*}
The trunk is the factorization of 6: $[2, 3]$.  The compatible spines are:
\begin{align*}
  P_C &= [2, 3] \quad \text{(product: }6\text{)} \\
  P_S &= [2, 3, 5] \quad \text{(product: }30\text{)} \\
  P_H &= [2, 3, 2, 5] \quad \text{(product: }60\text{)}
\end{align*}
The geographic interpretation is elegant: bisect Oregon into 2 halves, then
trisect each half into 3 regions (6 trunk regions = 6 congressional districts).
The senate subdivides each trunk region into 5 senate districts ($30/6 = 5$),
and the house further bisects each senate district into 2 house districts
($60/30 = 2$).  Every house district is contained in a senate district, every
senate district is contained in a congressional district.  No boundary zones.

\paragraph{Nesting map.} Figure~\ref{fig:or-nest} (schematic) illustrates this
three-level hierarchy for Oregon.  The geographic cuts at each level are
determined by GeoSection's ratio-optimal bisection of the state's census graph.

\begin{figure}[h]
\centering
\fbox{\parbox{0.8\textwidth}{\centering\vspace{2em}
\textit{[Oregon NestSection map --- 6 trunk regions, 30 senate subdivisions,}\\
\textit{60 house leaves --- to be generated by redist pipeline]}
\vspace{2em}}}
\caption{Schematic NestSection hierarchy for Oregon (C=6, S=30, H=60).
Dashed lines show trunk (congressional) boundaries; solid thin lines show
senate boundaries; gray fills show individual house districts.}
\label{fig:or-nest}
\end{figure}

\subsection{Case Study 2: North Carolina --- Partial Mode 3}

North Carolina has $C = 14$, $S = 50$, $H = 120$.  Then:
\begin{align*}
  g &= \gcd(14, 50, 120) = 2 \\
  \sigma &= (1 - 2/14) \times 100 \approx 85.7\%
\end{align*}
The trunk is $[2]$: a single east-west bisection of the state into two halves.
The compatible spines are:
\begin{align*}
  P_C &= [2, 7] \quad \text{(product: }14\text{)} \\
  P_S &= [2, 5, 5] \quad \text{(product: }50\text{)} \\
  P_H &= [2, 3, 4, 5] \quad \text{(product: }120\text{)}
\end{align*}
Within each half, the congressional map recurses to produce 7 districts; the
senate map produces 25 districts; the house map produces 60 districts.  The
senate and house are far from exact multiples of 7, requiring Mode~3 nesting
with boundary tolerance.

\paragraph{Nesting violation.}  The theoretical mismatch is driven by the tail
factorizations $[7]$ (congressional) vs.\ $[5,5]$ (senate).  In each half,
the 7 congressional boundaries do not align with the $5 \times 5$ senate
recursion.  The population fraction in boundary zones is determined empirically
by the geographic census-tract distribution; the plan specification reserves
$\tau_\text{pop} \leq 0.04$ (4\% of the state's population).

\paragraph{NestSection map for NC.}
The NestSection algorithm for NC proceeds as follows:
\begin{enumerate}
  \item Apply GeoSection to bisect NC into 2 trunk halves.
  \item In the western half: apply GeoSection(7) for congressional, GeoSection(25)
        for senate, NestRefine(25, 60/25) for house.
  \item In the eastern half: same recursion.
  \item Boundary zones are the census tracts that straddle a congressional/senate
        boundary; they are assigned to the larger-overlap chamber and flagged.
\end{enumerate}
This guarantees $V_\text{NC} \leq 0.04$ while sharing the single east-west trunk.

\subsection{Case Study 3: Texas --- Structural Incompatibility}

Texas has $C = 38$, $S = 31$, $H = 150$.  Then:
\begin{align*}
  g &= \gcd(38, 31, 150) = 1 \\
  \sigma &= (1 - 1/31) \times 100 \approx 96.8\%
\end{align*}
The trunk degenerates to the empty sequence $\tau = []$ (the product of no
factors is 1): the three chambers share only the state boundary.  The
compatible spines are simply the independent prime factorizations:
\begin{align*}
  P_C &= [2, 19] \\
  P_S &= [31] \quad\text{(prime)} \\
  P_H &= [2, 3, 5, 5]
\end{align*}
No geographic structure can be shared before the individual chamber
recursions diverge.  NestSection in Mode~3 with TX would require a very
large boundary tolerance ($\tau_\text{pop} > 0.1$) to force any
cross-chamber alignment, effectively making the constraint vacuous.

The fundamental obstacle is the primality of the senate count ($S = 31$):
a prime senate chamber is structurally incompatible with any congressional
count that is not a multiple of 31.  Texas is an extreme case, but the
pattern appears in several large states where the senate chamber size was
set by historical accident without regard for multi-level divisibility.

\subsection{Aggregate Findings}

\paragraph{Strict compatibility predictors.}  The 11 strictly compatible
states share two features: (a)~most have $C \leq 2$, making $g = C = \min$
trivially by the fact that $g \leq C \leq \min$ when $C$ divides $S$ and $H$;
(b)~the multi-seat strictly compatible states (AL, OR) have senate and house
counts that are exact multiples of $C$.  Oregon is the most structurally
interesting: $C = 6$, $S = 5C$, $H = 10C$, a clean $1:5:10$ ratio.
Alabama has $C = 7$, $S = 5C$, $H = 15C$, a clean $1:5:15$ ratio.

\paragraph{Score distribution.}  Among the 39 incompatible states, scores
cluster around three plateaus: $\{50\%, 66.7\%, 75\%\}$ for states with
small $C$ (2--4 seats) and modest GCDs; $\{80\%$--$90\%\}$ for mid-size
states; and $\{90\%$--$97\%\}$ for large states where congressional counts
are prime or coprime to chamber sizes.

\paragraph{Proposed nestability threshold.}  A statutory requirement for strict
nesting (Mode 1) where $g = \min(C,S,H)$ would apply to 11 states.  A weaker
requirement of $g \geq 5$ (meaningful shared trunk of at least 5 regions)
would apply to only 2 states: Alabama ($g=7$) and Oregon ($g=6$).  A
threshold of $g \geq 3$ would additionally include Louisiana ($g=3$), Arizona
($g=3$), New Jersey ($g=4$), California ($g=4$), Florida ($g=4$), Iowa ($g=2$),
Mississippi ($g=2$), Montana ($g=2$), North Carolina ($g=2$), and several others.
We revisit this threshold in Section~\ref{sec:discussion}.
